\documentclass[11pt]{article}
\usepackage[a4paper,margin=25mm]{geometry}
\usepackage{amsmath,amssymb,amsthm}
\usepackage[hidelinks]{hyperref}
\emergencystretch=2em
\newtheorem{theorem}{Theorem}
\newtheorem{lemma}{Lemma}
\title{Dirichlet approximation rules out the proposed\\logarithmic lower bound\\
\large Disproof of Conjecture 00000009690}
\author{gaochengzhecpu}\date{}
\begin{document}\maketitle
\begin{abstract}
Keep the algebraic arguments fixed at $2$ and $3$, and use integer
coefficients. For every $C>0$ and every height cutoff $H_0$ there
exist $a>0$, $b\geq0$ with actual height $H=\max(3,a,b)\geq H_0$
such that
\[
 0<|a\log2-b\log3|<\frac{C}{(\log H)^2}.
\]
Thus the conjectured lower bound fails already at $\varepsilon=1$,
even allowing an arbitrary positive implied constant and an
arbitrary sufficiently-large-height condition. The nonvanishing,
integer coefficient bounds and unbounded actual heights are all
proved in Lean 4.
\end{abstract}

\section*{Height convention and the statement being refuted}
The standard coefficient-height parameter for integer linear
forms is $\max(3,|b_1|,\ldots,|b_n|)$; see \cite[abstract]{bugeaud}.
We use $b_1=a$, $b_2=-b$ and fixed logarithms of the positive
rational algebraic numbers $2$ and $3$. Hence their heights and
degrees do not vary. Integer coefficients are a subclass of the
rational coefficients allowed by the conjecture.

The claimed improvement would in particular provide a positive
constant lower bound of order $(\log H)^{-2}$ for these fixed
arguments, once $H$ is sufficiently large. We disprove even
that version. This concerns the literal logarithmic lower bound
in the source, without claiming to determine any correct optimal
power $H^{-\kappa}$ in Baker's theory.

\begin{lemma}\label{lem:nonzero}
If $a>0$ and $b\geq0$ are integers, then
$a\log2-b\log3\ne0$.
\end{lemma}
\begin{proof}
Equality would imply $\log(2^a)=\log(3^b)$ and hence $2^a=3^b$.
The left side is divisible by $2$; the right side is not.
\end{proof}

\begin{lemma}\label{lem:approx}
For every integer $N\geq1$ there are integers $a,b$ with
$1\leq a\leq N$, $0\leq b\leq N+1$, and
\[
 0<|a\log2-b\log3|\leq\frac{\log3}{N+1}.
\]
\end{lemma}

\newpage
\begin{proof}
Put $\xi=\log2/\log3\in(0,1)$. Dirichlet's approximation
theorem gives $1\leq a\leq N$ and an integer $b$ with
$|a\xi-b|\leq1/(N+1)$. One proof considers the $N+1$ points
$0,\xi,\ldots,N\xi$ modulo $1$: one of their circular gaps
is at most $1/(N+1)$, and taking the two endpoint indices
gives the asserted approximation, including the possible
gap crossing $1$.

Because $a\xi\geq0$ and the error is strictly below $1$,
the integer $b$ is nonnegative. Also $a\xi\leq N$ and the
error is at most $1$, so $b\leq N+1$. Multiplication by
$\log3>0$ proves the upper bound; Lemma~\ref{lem:nonzero}
proves strict positivity.
\end{proof}

\begin{theorem}
For every $C>0$ and $H_0\in\mathbb N$ there are integers
$a>0$, $b\geq0$ such that $H=\max(3,a,b)\geq H_0$ and
\[
 0<|a\log2-b\log3|<C/(\log H)^2.
\]
\end{theorem}
\begin{proof}
The nonzero forms with $1\leq a\leq H_0$ and
$0\leq b\leq H_0$ form a finite set. There is therefore
$\delta>0$ below every absolute value in this set; if the
set is empty take $\delta=1$.

Choose $N\geq2$ large enough that
\[
 \frac{\log3}{N+1}<\delta,\qquad
 \frac{\log3\,(\log(N+1))^2}{N+1}<C.
\]
Both choices are possible by the elementary logarithmic
growth limit. Obtain $a,b$ from Lemma~\ref{lem:approx}.
Their actual height satisfies $H\leq N+1$. If $H<H_0$,
both coefficients would be at most $H_0$, contradicting
$|a\log2-b\log3|<\delta$. Thus $H\geq H_0$.
Since $H\geq3$, logarithms are positive and monotone, and
\[
 |a\log2-b\log3|(\log H)^2
 \leq \frac{\log3\,(\log(N+1))^2}{N+1}<C.
\]
Divide by $(\log H)^2>0$.
\end{proof}

This proves failure for every positive constant and every
cutoff at $\varepsilon=1$. A claim for every
$\varepsilon>0$ is consequently false. The actual maximum
height is used throughout; inflating a declared upper bound
is not the source of the contradiction.

\newpage
\section*{Formalization and reproducibility}
The Lean project uses the actual real logarithm and integer
coefficients. The nonvanishing lemma uses injectivity of
the logarithm on positive reals and prime divisibility.
The approximation lemma invokes Mathlib's proved Dirichlet
theorem and establishes coefficient bounds and sign conditions.

\texttt{logarithmic\_error\_limit} and
\texttt{reciprocal\_error\_limit} prove the two needed
limits. A finite-set induction proves a positive lower
bound for the finitely many small-height forms. The final
\texttt{counterexample} has arbitrary $C>0$ and $H_0$ and
uses \texttt{height a b = max 3 (max a b)}. The theorem
\texttt{no\_uniform\_bound} directly negates the existence
of a positive constant and a large-height cutoff for the
claimed $\varepsilon=1$ lower bound.

Lean 4.19.0 and Mathlib revision\\
\texttt{c44e0c8ee63ca166450922a373c7409c5d26b00b}
are pinned. From the submitted \texttt{lean/} directory:
\begin{verbatim}
lake exe cache get
lake build
lake env lean -DwarningAsError=true Main.lean
\end{verbatim}
The project is freshly built. The axiom audit contains only
standard foundational axioms, with no proof placeholders,
custom axioms or external computation oracles.

The optional exact-check script \texttt{verify.py} uses
only Python's standard library. It constructs rational
intervals enclosing logarithms from the convergent series
\[
 \log u=2\sum_{j\geq0}\frac{t^{2j+1}}{2j+1},
 \qquad t=\frac{u-1}{u+1},\quad 1\leq u\leq2,
\]
with a geometric remainder bound. It uses rational
continued fractions to find examples, then checks their
nonzero signs and error inequalities with exact fractions.
Floating point is used only to display the already bounded
errors. Finite examples do not prove the universal statement
or a limit; those are proved by Lean. Run
\texttt{python verify.py} and \texttt{tectonic main.tex}.

The exact bilingual conjecture is in \texttt{SOURCE.md};
provenance, real build logs, file hashes and PDF review are
in \texttt{verification/}. The author performed an adversarial
self-review. No independent reviewer or subagent was used.
The argument is an application of classical approximation,
not a claimed new result about linear forms in logarithms.

\begin{thebibliography}{9}
\bibitem{source} The Last Math Competition,
Conjecture 00000009690. Exact source and upstream provenance
are included with this submission.
\bibitem{bugeaud} Y. Bugeaud, \emph{$B'$}, 2022,
\href{https://arxiv.org/abs/2209.00275}{arXiv:2209.00275},
abstract (coefficient-height convention and background).
\bibitem{mathlib} The Mathlib community,
\emph{Diophantine Approximation}, Dirichlet approximation
theorem at the commit pinned above.
\href{https://github.com/leanprover-community/mathlib4/blob/c44e0c8ee63ca166450922a373c7409c5d26b00b/Mathlib/NumberTheory/DiophantineApproximation/Basic.lean}{Lean source}.
\end{thebibliography}
\end{document}
